\chapter{Multiplier ideals and test ideals}
\index{multiplier ideal}\index{test ideal}\index{asymptotic ideal}
\label{app:ideals}

Multiplier ideals and test ideals play parallel roles in the envelope
argument, but they are not the same construction.  This appendix maintains
three separate layers:
\begin{enumerate}
\item analytic multiplier ideals of quasi-plurisubharmonic functions on a
      complex manifold;
\item algebraic and asymptotic multiplier ideals on a smooth
      characteristic-zero variety;
\item test ideals and asymptotic test ideals on a regular \(F\)-finite
      scheme of characteristic \(p>0\).
\end{enumerate}
Only after each theory has been defined do we list the formal properties
used in \Cref{chap:bj-smooth-characteristic}.  No
characteristic-free identity between these sheaves is asserted.

\section{Complex analytic multiplier ideals}
\label{sec:appC-analytic-multiplier}

Let \(X\) be a complex manifold and let \(\varphi\) be quasi-psh.  Thus
locally \(\varphi\) is the sum of a psh function and a smooth function.

\begin{definition}[Analytic multiplier ideal]
\label{def:appC-analytic-multiplier}
The analytic multiplier ideal \(\mathcal J(\varphi)\) is the sheaf whose
stalk at \(x\in X\) consists of germs \(f\in\cO_{X,x}\) such that
\[
 |f|^2\rme^{-\varphi}
\tag{C.1}\label{eq:appC-analytic-integrability}
\]
is locally integrable near \(x\), with respect to one, hence any, smooth
positive volume form.
\end{definition}

The normalization in \eqref{eq:appC-analytic-integrability} is deliberate.
If
\[
 \varphi=c\log\left(\sum_{j=1}^r|g_j|^2\right)+O(1),
\tag{C.2}\label{eq:appC-analytic-singularity}
\]
then \(\mathcal J(\varphi)\) agrees with the algebraic convention
\(\mathcal J(\mathfrak a^c)\) for
\(\mathfrak a=(g_1,\ldots,g_r)\).  Authors using
\(|f|^2\rme^{-2\psi}\) write \(\psi=\varphi/2\).

\begin{theorem}\label{thm:appC-coherence-openness}
The sheaf \(\mathcal J(\varphi)\) is a coherent ideal.  It depends only on
the singularity type modulo locally bounded functions, and
\(\varphi\leq\psi+O(1)\) implies
\(\mathcal J(\varphi)\subset\mathcal J(\psi)\).  Moreover,
\[
 \mathcal J(\varphi)
   =\bigcup_{\varepsilon>0}\mathcal J((1+\varepsilon)\varphi)
\tag{C.3}\label{eq:appC-strong-openness}
\]
on relatively compact open sets, after adding a constant so that
\(\varphi\leq0\).
\end{theorem}

\begin{proof}[Elementary deductions and foundational inputs]
The bounded-invariance, monotonicity, and ideal properties follow directly
from local integrability.  We import coherence as Nadel's coherence theorem
in the normalization of \cite[Theorem~10.20]{Bou17}.  We import
\eqref{eq:appC-strong-openness} exactly from the strong openness theorem
\cite[Theorem~1.1]{GZ15}.  These analytic theorems are not inputs to the
algebraic envelope proof.
\end{proof}

For orientation, on a disc one computes
\(\mathcal J(c\log|z|^2)=(z^{\lfloor c\rfloor})\): if
\(f=z^m u\), \(u(0)\ne0\), integrability is exactly \(m>c-1\).

\section{Algebraic multiplier ideals}
\label{sec:appC-algebraic-multiplier}

Let \(X\) be smooth over a field of characteristic zero, and let
\(\mathfrak a\subset\cO_X\) be a nonzero coherent ideal.
Choose a log resolution
\[
 \mu:Y\longrightarrow X,\qquad
 \mathfrak a\cO_Y=\cO_Y(-F).
\tag{C.4}\label{eq:appC-log-resolution}
\]

\begin{definition}[Algebraic multiplier ideal]
\label{def:appC-algebraic-multiplier}
For \(c\in\mathbf R_{\geq0}\), define
\[
 \mathcal J_X(\mathfrak a^c)
 :=\mu_*\cO_Y\bigl(K_{Y/X}-\lfloor cF\rfloor\bigr).
\tag{C.5}\label{eq:appC-multiplier-resolution}
\]
\end{definition}

\begin{theorem}\label{thm:appC-independent-resolution}
The sheaf in \eqref{eq:appC-multiplier-resolution} is a coherent ideal of
\(\cO_X\) and is independent of the chosen log resolution.
\end{theorem}

\begin{proof}[From the local-vanishing input]
Given two log resolutions, resolve the indeterminacy of their birational
map and dominate both by a third.  For a composition
\(Z\xrightarrow{\nu}Y\xrightarrow{\mu}X\), the transformation rule follows
from
\[
 K_{Z/X}=K_{Z/Y}+\nu^*K_{Y/X}
\]
and local vanishing
\[
 R^i\nu_*\cO_Z
  \bigl(K_{Z/Y}-\lfloor\nu^*(cF)\rfloor\bigr)=0
 \quad(i>0).
\]
Projection formula then identifies the pushforwards.  Local vanishing is
the foundational input, in the form of
\cite[Theorem~9.4.1]{PAG2}.  Coherence follows from
properness of \(\mu\).
\end{proof}

For a prime divisor \(E\) on a smooth model over \(X\), write
\[
 A_X(\ord_E):=1+\ord_E(K_{Y/X}).
\tag{C.6}\label{eq:appC-log-discrepancy}
\]
Extend both sides homogeneously to positive multiples of divisorial
valuations.

\begin{theorem}\label{thm:appC-valuative-multiplier}
A regular function \(f\) on an open set \(U\subset X\) belongs to
\(\mathcal J_X(\mathfrak a^c)(U)\) if and only if
\[
 v(f)+A_X(v)>c\,v(\mathfrak a)
\tag{C.7}\label{eq:appC-valuative-multiplier}
\]
for every divisorial valuation \(v\) centered on \(U\).
\end{theorem}

\begin{proof}
On a fixed log resolution, membership in
\eqref{eq:appC-multiplier-resolution} says for each prime divisor \(E\)
\[
 \ord_E(f)+\ord_E(K_{Y/X})
       -\lfloor c\,\ord_E(F)\rfloor\geq0.
\]
Since the first two terms are integers, this is equivalent to
\[
 \ord_E(f)+A_X(\ord_E)>c\,\ord_E(\mathfrak a).
\]
A higher model gives the same criterion by resolution independence.
Every divisorial valuation appears on some higher smooth model.  This is
the algebraic valuative description; compare \cite{JM12}.
\end{proof}

\begin{corollary}\label{cor:appC-multiplier-valuation-bound}
For every divisorial \(v\),
\[
 v\bigl(\mathcal J_X(\mathfrak a^c)\bigr)
 \geq c\,v(\mathfrak a)-A_X(v).
\tag{C.8}\label{eq:appC-multiplier-valuation-bound}
\]
If \(\mathfrak a\subset\mathcal J_X(\mathfrak a)\), then also
\[
 v\bigl(\mathcal J_X(\mathfrak a)\bigr)\leq v(\mathfrak a).
\]
\end{corollary}

\begin{proof}
Apply \eqref{eq:appC-valuative-multiplier} to local generators of the
multiplier ideal and weaken the strict inequality.  The second assertion
uses reversal of order under ideal values.
\end{proof}

\begin{example}[Simple normal crossings]
\label{ex:appC-snc}
If \(X\) is smooth and
\(\mathfrak a=\cO_X(-F)\) with
\[
 F=\sum_i a_iD_i
\]
a simple normal crossings divisor, then the identity is a log resolution
and
\[
 \mathcal J_X(\mathfrak a^c)
 =\cO_X\left(-\sum_i\lfloor ca_i\rfloor D_i\right).
\tag{C.9}\label{eq:appC-snc-formula}
\]
This example checks both the floor and the strict inequality in
\eqref{eq:appC-valuative-multiplier}.
\end{example}

\section{Subadditivity and Skoda division}
\label{sec:appC-subadditivity-skoda}

\begin{theorem}\label{thm:appC-multiplier-subadditivity}
Let \(X\) be smooth in characteristic zero.  For nonzero ideals
\(\mathfrak a,\mathfrak b\) and \(s,t\geq0\),
\[
 \mathcal J_X(\mathfrak a^s\mathfrak b^t)
 \subset
 \mathcal J_X(\mathfrak a^s)\,
 \mathcal J_X(\mathfrak b^t).
\tag{C.10}\label{eq:appC-multiplier-subadditivity}
\]
In particular,
\[
 \mathcal J_X(\mathfrak a^{\ell t})
 \subset\mathcal J_X(\mathfrak a^t)^\ell
\qquad(\ell\in\mathbf N).
\tag{C.11}\label{eq:appC-multiplier-powers}
\]
\end{theorem}

\begin{proof}[Foundational input]
This is the subadditivity theorem of Demailly, Ein, and Lazarsfeld
\cite[Theorem~2.3]{DEL00}.  Its proof resolves the product ideal on
\(X^\ell\), applies the
restriction theorem to the diagonal, and uses the external-product formula
for multiplier ideals.  Smoothness of the diagonal is essential.
\end{proof}

\begin{theorem}\label{thm:appC-skoda-multiplier}
Let \(X\) be smooth of dimension \(n\).  For every nonzero coherent ideal
\(\mathfrak a\) and every real \(c\geq n\),
\[
 \mathcal J_X(\mathfrak a^c)
 =\mathfrak a\,
   \mathcal J_X(\mathfrak a^{c-1}).
\tag{C.12}\label{eq:appC-skoda-multiplier}
\]
\end{theorem}

\begin{proof}[From local vanishing]
Pull generators of \(\mathfrak a\) to a log resolution, where they have
the common factor defining \(F\).  The Koszul complex of a local reduction
of \(\mathfrak a\) with at most \(n\) generators, twisted by the divisors in
\eqref{eq:appC-multiplier-resolution}, is exact.  Local vanishing kills
higher direct images, so its first map gives
\eqref{eq:appC-skoda-multiplier}.  This is the geometric Skoda theorem;
see \cite[Section~9.6]{PAG2}.
\end{proof}

\begin{remark}
Skoda division controls high exponents by finitely many initial multiplier
ideals.  It is not used to identify multiplier ideals with integral
closures: the two operations have different discrepancy terms.
\end{remark}

\section{Asymptotic multiplier ideals}
\label{sec:appC-asymptotic-multiplier}

\begin{definition}[Graded sequence]
\label{def:appC-graded-sequence}
A graded sequence of ideals is a family
\(\mathfrak a_\bullet=(\mathfrak a_m)_{m\geq1}\) with
\[
 \mathfrak a_m\mathfrak a_\ell
 \subset\mathfrak a_{m+\ell}.
\tag{C.13}\label{eq:appC-graded-sequence}
\]
It is nonzero if some, hence all sufficiently divisible, terms are nonzero.
For a valuation \(v\), subadditivity of the numerical sequence
\(v(\mathfrak a_m)\) gives
\[
 v(\mathfrak a_\bullet)
 :=\inf_{m\geq1}\frac1m v(\mathfrak a_m)
 =\lim_{\substack{m\to\infty\\ \mathfrak a_m\ne0}}
       \frac1m v(\mathfrak a_m)
\tag{C.14}\label{eq:appC-asymptotic-order}
\]
where the limit is taken only through the nonzero terms.  Equivalently,
it is the limit along every cofinal divisibility net of nonzero indices.
This qualification is necessary because ``nonzero'' for a graded
sequence means that one term is nonzero; it does not rule out zero terms
at other indices.
\end{definition}

\begin{definition}[Asymptotic multiplier ideal]
\label{def:appC-asymptotic-multiplier}
For \(t\geq0\), the ideals
\[
 \mathcal J_X\left(\mathfrak a_q^{\,t/q}\right),
\qquad q\geq1,\quad \mathfrak a_q\ne0,
\]
have a unique maximal member.  It is denoted
\[
 \mathcal J_X(\mathfrak a_\bullet^t).
\tag{C.15}\label{eq:appC-asymptotic-multiplier}
\]
For sufficiently divisible nonzero indices \(q\),
\[
 \mathcal J_X(\mathfrak a_\bullet^t)
 =\mathcal J_X(\mathfrak a_q^{\,t/q}).
\tag{C.16}\label{eq:appC-asymptotic-stabilization}
\]
\end{definition}

\begin{proof}[Why a maximal member exists]
If \(q\mid r\), gradedness gives
\(\mathfrak a_q^{r/q}\subset\mathfrak a_r\).  Monotonicity of multiplier
ideals then compares the corresponding members in the direction needed to
form an increasing divisibility net.  Since \(X\) is noetherian, coherent
ideals satisfy the ascending-chain condition, so the net stabilizes.  Any
two indices have a common multiple, giving uniqueness of the maximum.
\end{proof}

\begin{proposition}\label{prop:appC-asymptotic-multiplier-properties}
For a nonzero graded sequence on a smooth characteristic-zero variety:
\begin{enumerate}
\item \(\mathfrak a_m\subset
      \mathcal J_X(\mathfrak a_\bullet^m)\);
\item for \(m,\ell\geq1\),
\[
 \mathcal J_X(\mathfrak a_\bullet^{m\ell})
 \subset
 \mathcal J_X(\mathfrak a_\bullet^m)^\ell;
\tag{C.17}\label{eq:appC-asymptotic-subadditivity}
\]
\item for every divisorial \(v\),
\[
 v\bigl(\mathcal J_X(\mathfrak a_\bullet^m)\bigr)
 \geq m\,v(\mathfrak a_\bullet)-A_X(v).
\tag{C.18}\label{eq:appC-asymptotic-discrepancy}
\]
\end{enumerate}
\par
\end{proposition}

\begin{proof}
Choose \(q\) divisible by \(m\) and sufficiently divisible for
\eqref{eq:appC-asymptotic-stabilization}.  Gradedness gives
\(\mathfrak a_m^{q/m}\subset\mathfrak a_q\).  Thus, for
\(f\in\mathfrak a_m\) and every divisorial \(v\),
\[
 v(f)\geq v(\mathfrak a_m)
       \geq\frac mqv(\mathfrak a_q).
\]
Since \(A_X(v)>0\) on smooth \(X\), the strict criterion
\eqref{eq:appC-valuative-multiplier} puts \(f\) in
\(\mathcal J(\mathfrak a_q^{m/q})\), proving (1).  Apply
\Cref{thm:appC-multiplier-subadditivity} before stabilization to get
(2).  Applying \eqref{eq:appC-multiplier-valuation-bound} to
\(\mathfrak a_q^{m/q}\) gives
\[
 v\bigl(\mathcal J(\mathfrak a_\bullet^m)\bigr)
 \geq\frac mqv(\mathfrak a_q)-A_X(v)
 \geq m\,v(\mathfrak a_\bullet)-A_X(v),
\]
which is (3).  See \cite[Chapter~11]{PAG2} and \cite{JM12}.
\end{proof}

\begin{remark}[Valuative safeguard]
\label{rem:appC-asymptotic-valuative-safeguard}
For a fixed ordinary ideal, \Cref{thm:appC-valuative-multiplier} is
an exact if-and-only-if.  For a graded sequence, the definition used in
this book is the stabilized maximum
\eqref{eq:appC-asymptotic-multiplier}; the envelope proof needs only the
one-sided estimate \eqref{eq:appC-asymptotic-discrepancy}.  We do not
replace stabilization by an unproved pointwise criterion over arbitrary
semivaluations.
\end{remark}

\section{Uniform global generation}
\label{sec:appC-global-generation}

The analytic or valuative estimate identifies the limiting ideal function,
but psh regularization also requires globally generated approximants.

\begin{theorem}\label{thm:appC-uniform-generation-multiplier}
Let \(X\) be smooth projective of characteristic zero, let \(L\) be a line
bundle, and let \(\mathfrak a_\bullet\) be a graded sequence such that
\(\mathfrak a_q(qL)\) is globally generated for all sufficiently divisible
\(q\).  Then there exists a line bundle \(A\), independent of \(m\), such
that
\[
 \mathcal J_X(\mathfrak a_\bullet^m)(mL+A)
\tag{C.19}\label{eq:appC-uniform-generation-multiplier}
\]
is globally generated for every sufficiently divisible \(m\).
\end{theorem}

\begin{proof}[Foundational uniform-generation input]
Choose a sufficiently positive very ample bundle \(H\).  Stabilize
\(\mathcal J(\mathfrak a_\bullet^m)\) at a divisible index \(q\), resolve
\(\mathfrak a_q\), and apply local vanishing on the resolution.
After twisting by \(mL+A-iH\), \(1\leq i\leq n\), the positivity supplied
by the global generation of \(\mathfrak a_q(qL)\) gives the vanishings
required by Castelnuovo--Mumford regularity.  The same fixed \(A\) works
for all \(m\).  This is the uniform generation input used in the proof of
\cite[Theorem~5.20]{BJ22}; see also \cite[Chapter~11]{PAG2}.
\end{proof}

\begin{remark}[Checkpoint]
The fixed twist in \eqref{eq:appC-uniform-generation-multiplier} is
essential.  If the auxiliary line bundle were allowed to grow linearly
with \(m\), the normalized numerical classes would not converge to
\(c_1(L)\).
\end{remark}

\section{Test ideals in positive characteristic}
\label{sec:appC-test-ideals}

Let \(S\) be a regular \(F\)-finite scheme of characteristic \(p>0\).
\(F\)-finiteness makes Frobenius finite, while regularity makes it flat by
Kunz's criterion.  In particular, its pushforwards of coherent sheaves are
coherent.

\begin{definition}[Generalized test ideal]
\label{def:appC-test-ideal}
For an ideal \(J\subset\cO_S\), write \(J^{[1/p^e]}\) for the smallest
ideal \(K\) such that \(J\subset K^{[p^e]}\).  For a nonzero ideal
\(I\subset\cO_S\) and \(t\geq0\), the generalized test ideal is the stable
value, computed locally, of the increasing sequence
\[
 \left(I^{\lceil tp^e\rceil}\right)^{[1/p^e]}.
\]
If \(S\) admits an invertible dualizing sheaf \(\omega_S\), this local
definition is equivalently characterized, after tensoring by
\(\omega_S\), by the stabilized Frobenius-trace sum
\[
 \tau_S(I^t)\otimes\omega_S
 =
 \sum_{e\geq0}
 \operatorname{Tr}^e\left(
 F_*^e(I^{\lceil tp^e\rceil}\otimes\omega_S)
 \right).
\tag{C.20}\label{eq:appC-test-trace}
\]
\end{definition}

\begin{remark}[Conventions]
The Frobenius-root ideals, and equivalently the trace images when
\(\omega_S\) is available, form an increasing sequence and stabilize by
noetherianity.  The schemes \(S=X\times\mathbf A^1\) used below possess
the required invertible dualizing sheaf.
Equivalent definitions sometimes place \(p^e-1\) rather than \(p^e\) in
an exponent, especially for pairs and Cartier algebras.  Formula
\eqref{eq:appC-test-trace} is the convention used by the envelope proof.
Changing conventions requires tracking the rounding error; it is never
silently discarded.
\end{remark}

\begin{definition}[Asymptotic test ideal]
\label{def:appC-asymptotic-test}
For a nonzero graded sequence \(\mathfrak a_\bullet\), define
\[
 \tau_S(\mathfrak a_\bullet^t)
 :=\max_{\mathfrak a_q\ne0}\tau_S(\mathfrak a_q^{t/q}).
\tag{C.21}\label{eq:appC-asymptotic-test}
\]
Noetherian stabilization gives, for sufficiently divisible nonzero
indices \(q\),
\[
 \tau_S(\mathfrak a_\bullet^t)
 =\tau_S(\mathfrak a_q^{t/q}).
\tag{C.22}\label{eq:appC-test-stabilization}
\]
\end{definition}

\begin{theorem}\label{thm:appC-test-properties}
For integers \(m,\ell\geq1\),
\[
 \mathfrak a_m\subset\tau_S(\mathfrak a_\bullet^m),
\qquad
 \tau_S(\mathfrak a_\bullet^{m\ell})
 \subset\tau_S(\mathfrak a_\bullet^m)^\ell.
\tag{C.23}\label{eq:appC-test-subadditivity}
\]
If \(w\) is divisorial and centered on \(S\), then
\[
 w(\tau_S(I^t))
 \geq t\,w(I)-A_S(w).
\tag{C.24}\label{eq:appC-test-discrepancy}
\]
\end{theorem}

\begin{proof}[Foundational inputs and deduction]
The containment and asymptotic subadditivity are
\cite[Section~3 and Proposition~3.1]{MustataTest}.  For
\eqref{eq:appC-test-discrepancy}, trivialize \(\omega_S\) near the center
of \(w\).  The trace becomes a \(p^{-e}\)-linear map
\(\Phi_e:F_*^e\cO_S\to\cO_S\).  If \(g=\Phi_e(f)\ne0\), Canton's
Cartier-algebra discrepancy gives
\[
 \frac{w(f)-p^ew(g)}{p^e-1}\leq A_S(w).
\]
Taking \(f\in I^{\lceil tp^e\rceil}\), minimizing over trace generators,
and stabilizing proves \eqref{eq:appC-test-discrepancy}; see
\cite[Definition~3.2 and Proposition~4.1]{Canton}.
\end{proof}

\begin{theorem}\label{thm:appC-skoda-test}
If \(I\) is locally generated by at most \(r\) elements, then
\[
 \tau_S(I^t)=I\,\tau_S(I^{t-1})
 \qquad(t\geq r).
\tag{C.25}\label{eq:appC-skoda-test}
\]
\end{theorem}

\begin{proof}[Foundational input]
This is the Brian\c con--Skoda property for generalized test ideals; see
\cite[Theorem~6.9(iii)]{ST12}; see also
\cite[Theorem~2.1]{HY03}.  These sources state the rational-exponent form.
The real-exponent form follows by choosing a decreasing rational sequence
on which the test ideal is constant, using right-continuity
\cite[Theorem~6.6]{ST12}.  The proof expands sufficiently high powers of an
\(r\)-generated ideal and passes the identity through Frobenius roots.
\end{proof}

\begin{proposition}\label{prop:appC-test-fixed-twist}
Let \(X\) be smooth projective of dimension \(n\) over an \(F\)-finite
field, put \(S=X\times\mathbf A^1\), and let \(J_\bullet\) be the graded
base-ideal sequence associated to a line bundle \(L\) in
\Cref{chap:bj-smooth-characteristic}.  Write again \(L\) for its pullback
to \(S\).
Choose a very ample globally generated \(H\), \(d\geq n+1\), and set
\[
 A=K_X+dH.
\]
Write also \(H\) and \(A\) for their pullbacks to \(S\).
Then
\[
 \tau_S(J_\bullet^m)(mL+A)
\tag{C.26}\label{eq:appC-test-fixed-twist}
\]
is generated relative to \(\mathbf A^1\) for every \(m>0\).
\end{proposition}

\begin{proof}
This is \Cref{prop:bj-smooth-uniform-generation}.  Its proof
uses the quotient supplied by \eqref{eq:appC-test-trace}, the projection
formula for finite Frobenius, and relative
Castelnuovo--Mumford regularity.  The residual coefficient
\[
 mp^e-q\left\lceil\frac{mp^e}{q}\right\rceil
\]
lies in a fixed finite interval, while the coefficient
\(p^e(d-i)\) of \(H\) tends to \(+\infty\) for \(1\leq i\leq n\).
This is the exact reason a twist independent of \(m\) is available.
\end{proof}

\section{Comparison across characteristic}
\label{sec:appC-comparison-characteristic}

\begin{theorem}\label{thm:appC-hara-yoshida}
Let \(X,\mathfrak a,c\) be defined in characteristic zero, with
\(\mathfrak a\) an algebraic ideal and \(c\in\mathbf Q_{\geq0}\).
After spreading the data over a finitely generated
\(\mathbf Z\)-algebra, for closed fibres in a suitable dense open set one
has
\[
 \mathcal J_X(\mathfrak a^c)_p
 =\tau_{X_p}(\mathfrak a_p^c).
\tag{C.27}\label{eq:appC-reduction-comparison}
\]
\end{theorem}

\begin{proof}[Foundational input and scope]
This is \cite[Theorem~6.8]{HY03}.  It concerns reductions of one fixed datum
after excluding primes; it neither defines multiplier ideals in
characteristic \(p\) nor supplies uniform equality for arbitrary graded
sequences.
\end{proof}

\begin{remark}[The characteristic firewall]
\label{check:appC-characteristic-firewall}
Analytic multiplier ideals use integrability; algebraic multiplier ideals
use characteristic-zero log resolutions; test ideals use Frobenius trace.
Equality \eqref{eq:appC-reduction-comparison} holds only after spreading
out and excluding primes.  The envelope proof runs the two algebraic
engines separately and never treats this as a characteristic-free identity.
\end{remark}

\section{The exact parallel in the envelope proof}
\label{sec:appC-envelope-parallel}

Let \(X\) be smooth projective over the trivially valued ground field and
let
\[
 S=X\times\mathbf A^1.
\]
After resolving a rational PL obstacle, \cref{chap:bj-smooth-characteristic} constructs a graded
sequence \(J_\bullet\) on \(S\) satisfying
\[
 J_mJ_r\subset J_{m+r}.
\tag{C.28}\label{eq:appC-envelope-graded-sequence}
\]
For \(m>0\), define
\[
 B_m=
 \begin{cases}
 \mathcal J_S(J_\bullet^m),
      &\operatorname{char}k=0,\\
 \tau_S(J_\bullet^m),
      &\operatorname{char}k=p>0.
 \end{cases}
\tag{C.29}\label{eq:appC-envelope-two-engines}
\]
The proof imports exactly the following four statements.

\begin{center}
\begin{tabular}{p{0.20\textwidth}p{0.31\textwidth}p{0.31\textwidth}}
\toprule
Required output & Multiplier ideal & Test ideal\\
\midrule
Containment
 & \(J_m\subset B_m\)
 & \(J_m\subset B_m\)\\
Subadditivity
 & \(B_{m\ell}\subset B_m^\ell\)
 & \(B_{m\ell}\subset B_m^\ell\)\\
Fixed twist
 & uniform generation theorem
 & Frobenius trace and
   \eqref{eq:appC-test-fixed-twist}\\
Valuation error
 & \(v(B_m)\geq m v(J_\bullet)-A_S(v)\)
 & the same inequality from Canton\\
\bottomrule
\end{tabular}
\end{center}

Writing \(\phi_I=-v(I)\), containment and subadditivity imply that
\[
 u_m:=m^{-1}\phi_{B_m}
\tag{C.30}\label{eq:appC-envelope-stages}
\]
decreases along divisibility and lies above the normalized base-ideal
limit.  Fixed-twist generation makes \(u_m\) a GF psh function in the
class \(L+A/m\).  For each fixed divisorial valuation, the discrepancy
error divided by \(m\) tends to zero.  Thus the limiting psh function has
the correct divisorial values.

\begin{remark}[Supported points]
\label{rem:appC-supported-points}
The ideal-theoretic valuation estimates are made at full-support
divisorial valuations, where all ideal values are finite.  At a
semivaluation whose support contains an ideal, \(\phi_I=-\infty\).
\Cref{chap:bj-smooth-characteristic} retains those values and uses
global psh comparison to pass from divisorial equality to equality on the
whole analytification.  It never subtracts two infinite values.
\end{remark}

\begin{remark}[What is not imported]
No step requires resolution in characteristic \(p\), continuity of every
ideal function, or a valuative criterion over supported semivaluations.
Skoda division is structural background; it does not replace fixed-twist
generation or the discrepancy estimate.  Finally, subadditivity becomes
monotonicity only after applying the sign-reversing operation \(I\mapsto
-v(I)\), and all divisibility limits are nets.
\end{remark}

\section{Dependency map}
\label{sec:appC-dependency-map}

\begin{center}
\begin{tabular}{p{0.34\textwidth}p{0.51\textwidth}}
\toprule
Ideal-theoretic tool & First substantial use\\
\midrule
Analytic multiplier ideals
 & singularity theory in \cref{part:complex}\\
Log-resolution formula
 & characteristic-zero regularization,
   \Cref{chap:bj-smooth-characteristic}\\
Asymptotic subadditivity
 & decreasing ideal approximants in the same chapter\\
Uniform global generation
 & convergence of numerical classes in the same chapter\\
Test-ideal trace and subadditivity
 & positive-characteristic smooth case in the same chapter\\
Canton discrepancy estimate
 & divisorial identification in the same chapter\\
Reduction modulo \(p\)
 & comparison only; not a dependency of the envelope proof\\
\bottomrule
\end{tabular}
\end{center}

\section{Source note}
Multiplier ideals, local vanishing, Skoda division, and asymptotic ideals
are treated in \cite{PAG2}; subadditivity is \cite{DEL00}, and valuation
asymptotics are developed in \cite{JM12}.  Generalized test ideals and
their multiplier-ideal comparison originate in \cite{HY03}; see
\cite{ST12} for conventions.  The particular asymptotic test-ideal
properties used here are from \cite{MustataTest}, and the discrepancy
estimate is from \cite{Canton}.  The non-Archimedean envelope use is
\cite[Theorem~5.20]{BJ22}, \cite{BJ22b}, and
\Cref{chap:bj-smooth-characteristic}.
